\subsection{Datasets}

\paragraph{MSLesSeg.}
The primary experiments use the MSLesSeg
dataset~\cite{guarnera_mslesseg_2025} of longitudinal brain MRI scans of
patients with multiple sclerosis. Of its 75 subjects, the 53 with manual
lesion annotations are used here, as the official test set is unlabeled.
Each subject has one to four timepoints, each with T1-WI, T2-WI, and Fluid-Attenuated Inversion Recovery (FLAIR)
images. The dataset organizers provide the images skull-stripped,
co-registered to the MNI152 template, and resampled to 1~mm isotropic
spacing.

\paragraph{3D-MR-MS.}
To assess whether the behavior of CATMIL carries over to other data, the
five losses are also trained and evaluated on a second multiple sclerosis
dataset, 3D-MR-MS~\cite{lesjak_novel_2018}, which was acquired at a
different center with different imaging conditions. It contains 30
patients with consensus-based white matter lesion annotations, provided
co-registered and bias-field corrected. Of the available modalities, only
T1-WI, T2-WI, and FLAIR are used, to keep the input configuration identical
to MSLesSeg. The voxels of this dataset are anisotropic
($0.47 \times 0.47 \times 0.8$~mm), so lesion sizes on 3D-MR-MS are given
in mm$^3$ rather than in voxels.

\subsection{Data Split}

Data are partitioned at the patient level, so that scans of one patient
never appear in both training and evaluation sets. For MSLesSeg, six
patients, with all their timepoints (12 test cases), are held out as an
independent test set; for 3D-MR-MS, six of the 30 patients are held out.
The remaining patients (47 and 24, respectively) are divided into five
folds following the nnU-Net cross-validation strategy. Each of the five
fold models is evaluated on the test set, and results are reported as the
mean and sample standard deviation across the five models; statistical
comparisons are described in \cref{sec:stats}. All training
settings, including the CATMIL loss weights, are identical for the two
datasets, so 3D-MR-MS was not used to tune the objective.

\subsection{Compared Losses}

Five configurations of nnU-Net are compared, differing only in the loss
function:
\begin{itemize}
    \item \textbf{DiceCE}: the default Dice and cross-entropy loss;
    \item \textbf{Tversky}~\cite{salehi_tversky_2017}: a loss that
    penalizes missed lesion voxels more than false-positive voxels
    ($\alpha=0.3$, $\beta=0.7$);
    \item \textbf{FocalTversky}~\cite{abraham_novel_2019}: a Tversky loss
    that additionally emphasizes difficult voxels (focal exponent
    $\gamma_{\mathrm{F}}=4/3$);
    \item \textbf{CC-DiceCE}~\cite{bouteille_learning_2026}: a published
    loss that adds a per-lesion DiceCE term, weighting each lesion equally,
    to the global DiceCE loss;
    \item \textbf{CATMIL}: the proposed loss.
\end{itemize}
CC-DiceCE is the closest prior objective: like the CAT term, it balances
lesions of different sizes, but it scores each lesion by overlap and has no
detection term. It was ported unchanged into the same nnU-Net pipeline and
150-epoch budget. Ablations of the two CATMIL terms and of their weights
are reported in
\cref{appendix:catmil-components,appendix:catmil-loss-weights}.

\subsection{Implementation}

All models use nnU-Net (version 2.1.1)~\cite{isensee_nnu-net_2021} with
its default preprocessing, patch-based training, and optimizer, trained for
150 epochs (details in \cref{appendix:impl-stats}, `Implementation'). The loss
implementation and evaluation code are available at
\url{https://github.com/luumsk/SmallLesionMRI}.

\subsection{Evaluation Metrics}

Voxel-level metrics describe overall segmentation quality, and lesion-level
metrics describe detection, in particular of small lesions. Lesion-level
metrics are needed because the DSC is dominated by large lesions
(\cref{sec:problem}) and does not reflect whether individual, clinically
relevant lesions are found~\cite{basit_rethinking_2026,jaus_every_2025}.

\paragraph{Voxel-level metrics.}
Overall segmentation quality is measured with the DSC~\cite{dice_measures_1945}
and the 95th-percentile Hausdorff distance (HD95, in
mm)~\cite{taha_metrics_2015}, a measure of boundary error.

\paragraph{Lesion definition and detection.}
A lesion is a 3D connected component of the binary mask under
6-connectivity, the same definition as in the proposed loss. A reference
lesion counts as detected if at least one predicted voxel overlaps it, and
a predicted component counts as correct if it overlaps at least one
reference lesion. This any-overlap criterion is lenient, so lesion recall
was also computed with stricter criteria that require at least 10\%, 25\%,
or 50\% of the lesion voxels to be covered
(\cref{appendix:strict-detection}).

\paragraph{Lesion-level metrics.}
Lesion-wise recall is the fraction of reference lesions that are detected,
lesion-wise precision is the fraction of predicted components that are
correct, and the lesion-wise F1 score combines the two. Because precision
counts components regardless of their size, many small false-positive
components lower it substantially even when their total volume is small.
Small Lesion Recall is the lesion-wise recall restricted to reference
lesions of at most 150 voxels, which corresponds to 150~mm$^3$ at the 1~mm
isotropic spacing of MSLesSeg and to about 26~mm$^3$ on 3D-MR-MS. The
metric therefore covers a much smaller physical size range on 3D-MR-MS, so
its absolute values are not comparable between the two datasets and are
used only to compare losses within a dataset; the comparison across
datasets relies on the recall of lesions of at least 3~mm, which is defined
in physical units. The FN
Count is the number of reference lesions that are completely missed. The
recall of lesions of at least 3~mm is the lesion-wise recall restricted to
lesions above the clinical size threshold (\cref{sec:intro};
\cite{barkhof_2024_2025,filippi_mri_2016}); the longest diameter of a
reference lesion is measured as its largest in-plane extent over all
axial, coronal, and sagittal slices.

\paragraph{Error-volume metrics.}
The FP Volume is the physical volume of the false-positive voxels, and the
FN Volume Fraction is the share of the total reference lesion volume that
belongs to missed lesions. All metrics are computed per test case and
averaged over cases; their exact definitions are given in
\cref{appendix:metrics}.

\subsection{Statistical Analysis}
\label{sec:stats}

The standard deviations in the result tables describe the variation
between the five fold models, not the uncertainty from the small test sets.
Differences between CATMIL and each baseline were therefore also assessed
with the patient as the independent unit, which matters for MSLesSeg, whose
12 test cases are longitudinal scans of six patients. For each metric, a
95\% confidence interval (CI) was obtained with a patient-level bootstrap
(10{,}000 resamples), and the number of patients in which CATMIL was better
is reported. With six patients, no Holm-corrected paired Wilcoxon test can
reach $p < 0.05$, so a difference is regarded as supported when its 95\%
CI excludes zero; these analyses describe the consistency of an effect
across patients, not population-level estimates. Details are given in
\cref{appendix:impl-stats}, `Statistical analysis'. To obtain a larger sample, the MSLesSeg comparisons
are also repeated on the cross-validation predictions of the 47 training
patients (\cref{sec:catmil_oof}).
